\begin{workedexample}[Worked Example: 3~dB Pi attenuator in 50~$\Omega$]
For attenuation $A$ (voltage ratio $K = 10^{A/20}$) in a $Z_0$ system, the
Pi-pad resistors are
\[ R_{\text{series}} = Z_0\frac{K^2-1}{2K}, \qquad
   R_{\text{shunt}} = Z_0\frac{K+1}{K-1}. \]
For $A = 3\ \text{dB}$, $K = 1.413$: $R_{\text{series}} = 17.6\ \Omega$ and
$R_{\text{shunt}} = 292\ \Omega$. The pad attenuates $3\ \text{dB}$ while
presenting $50\ \Omega$ at both ports.
\end{workedexample}

\begin{keytakeaways}
\begin{itemize}[leftmargin=1.4em,itemsep=2pt]
\item Resistive Pi/T pads attenuate while keeping both ports matched to $Z_0$.
\item Values follow from $K = 10^{A/20}$; symmetric pads are bidirectional.
\item The pad dissipates the lost power --- size the resistors for it.
\end{itemize}
\end{keytakeaways}

\begin{exercises}
\item What voltage ratio $K$ corresponds to a $6\ \text{dB}$ pad?
\item What fraction of input power reaches the load through a $10\ \text{dB}$ pad?
\item A $20\ \text{dB}$ pad passes $1\ \text{W}$ in. Roughly how much must it dissipate?
\end{exercises}

\furtherreading{Pozar~\cite{pozar}; Hayward~\cite{hayward}.}
